\documentclass[11pt]{article}
\usepackage[utf8]{inputenc}
\usepackage[T1]{fontenc}
\usepackage{geometry}
\geometry{a4paper, margin=1in}

\begin{document}

\begin{center}
\Large\bfseries Endorsement Request --- 校友视角
\end{center}

\vspace{0.3cm}

\noindent\textbf{From:} 20242100136@dcc.edu.cn \\
\textbf{To:} wxf@cjlu.edu.cn（中国计量大学 汪晓锋老师） \\
\textbf{Subject:} 滇池学院校友 arXiv cs.LG endorsement 请求 --- 跨域 CSI 手势识别可复现性研究

\vspace{0.3cm}
\hrule
\vspace{0.3cm}

\noindent 汪老师：

您好！冒昧来信。我是云南大学滇池学院电子信息工程专业的访问研究人员。了解到您是滇池学院 2009 届电子信息工程专业校友（同系前辈），并且长期在中国计量大学从事图神经网络、推荐系统、图异常检测、大模型技术等 AI 方向研究，发表 30+ 篇论文，特此冒昧请求您的 endorsement 支持。

\vspace{0.2cm}

\noindent 我最近准备向 arXiv 投稿一篇论文：

\begin{quote}
\textbf{标题：}WiFi CSI Cross-Domain Gesture Recognition: A Reproducibility Audit and the Boundary of Multi-Rx Representation

\textbf{类别：}cs.LG（primary）+ eess.SP（cross-list）

\textbf{核心发现：}在 18-fold × 5-seed paired 协议下，10 种跨域方法的"提升"全部不显著（$|d|<0.3$, $p>0.28$），核心原因是 seed-to-seed std（4.82 pp）远大于典型报告提升。多 Rx 表示在 Widar3.0 上 +8.14 pp 但在 MMFi 上不迁移——说明"没有普遍最优的 CSI 表示"。我们也实证检验了 Liu 2025 (arXiv 2512.04521) 的 CBAM+ResNet18 报告——在严格的 paired × multi-seed 协议下，标准 CBAM+ResNet18 跨域仅 28.35\%，远低于其报告的 97.61\%。
\end{quote}

\vspace{0.2cm}

\noindent\textbf{论文与您研究的相关性：}
\begin{itemize}
\item 您长期做图神经网络，与我们的网络架构分析直接相关
\item 您做推荐系统 + 大模型，对 ML benchmark 方法论有深度理解
\item 论文核心议题（跨域泛化、可复现性、benchmark 协议）属于 cs.LG 主线
\item 校友身份（同为滇池学院电子信息工程系背景），希望前辈支持
\end{itemize}

\vspace{0.2cm}

\noindent\textbf{endorsement 操作：}
\begin{enumerate}
\item 您登录自己的 arXiv 账号（您应该有 cs.LG/cs.AI 的 endorsement 资格）
\item 访问 https://arxiv.org/auth/endorsers
\item endorse others → 搜索我的 submission ID（投稿后补充给您）
\item 点击 endorse
\end{enumerate}

\vspace{0.2cm}

\noindent\textbf{补充材料：}
\begin{itemize}
\item \texttt{arxiv\_submission/paper\_arxiv.tex}（LaTeX 源文件）
\item \texttt{arxiv\_submission/paper\_arxiv.pdf}（编译后 PDF）
\item \texttt{arxiv\_submission/ARXIV\_METADATA.txt}（投稿元数据）
\end{itemize}

\vspace{0.2cm}

\noindent 如蒙支持，万分感谢。如您工作繁忙无法 endorse，也请告知，方便我及时调整投稿策略。

\vspace{0.3cm}

\noindent 此致 \\
敬礼

\vspace{0.3cm}

\noindent [你的名字] \\
20242100136@dcc.edu.cn \\
滇池学院（云南大学滇池学院） \\
电子信息工程专业

\end{document}
